\subsection{תהליך הלמידה}

אחד הדברים המרכזיים שלמדתי במהלך הפרויקט הוא שלמידת מכונה אמיתית אינה
נגמרת באימון מודל. בתחילת הדרך חשבתי שהאתגר המרכזי יהיה לבחור
ארכיטקטורה טובה, אבל תוך כדי העבודה הבנתי שחלק גדול מן הקושי נמצא
באיסוף הנתונים, בתיוג, ובהרכבת כל הרכיבים למערכת אחת שעובדת מקצה לקצה.

בנוסף, למדתי לעבוד עם מגוון טכנולוגיות --- \textenglish{YOLO},
\textenglish{ByteTrack}, \textenglish{FastAPI}, \textenglish{Flutter}
ועיבוד וידאו --- ולחבר ביניהן לפתרון אחד.

\subsection{הרגע הקשה ביותר}

הרגע הקשה ביותר מבחינתי היה להבין שזיהוי הספרות לא יפתור את כל הבעיה
לבד. היו שלבים שבהם היה נדמה שאם רק אצליח לשפר עוד קצת את מודל הספרות,
הכל יעבוד. אבל בפועל, גם עם מודל טוב יותר עדיין נשארו חסימות, טשטושי
תנועה, זוויות קשות והחלפות זהות.

ההבנה הזאת הייתה מתסכלת בהתחלה, אבל היא גם הכריחה אותי לשנות גישה:
במקום לשאול רק "איך לשפר את המודל", התחלתי לשאול "איך לגרום למערכת
כולה לקבל החלטות טובות יותר". משם הגיעו הרעיונות של שימוש בלוח
המשחקים, מנגנון הצבעה מצטבר ומרווח ביטחון דינמי.

\subsection{הרגע המרגש ביותר}

הרגע המרגש ביותר היה בפעם הראשונה שבה ראיתי את ה־\textenglish{Replay}
של המסלולים המחושבים, ובמיוחד כשהם התיישבו בצורה משכנעת עם הווידאו של
המשחק. זה היה הרגע שבו הרגשתי שהפרויקט עבר ממצב של ``קוד שעובד חלקית''
למשהו שבאמת מרגיש כמו מוצר קטן ושימושי.

גם לראות איך כמה רכיבים נפרדים --- זיהוי, מעקב, זיהוי ספרות,
\textenglish{Levenshtein} והחלקת מסלולים --- מתחברים לתוצאה אחת, היה
מאוד מספק.

\subsection{ניהול זמן ומה הייתי עושה אחרת}

אם יש משהו שהייתי עושה אחרת, זה לבנות מוקדם יותר פרוטוקול הערכה מסודר
ולשמור תיעוד מדויק יותר של כל ריצת אימון. במהלך הפרויקט היו הרבה
ניסויים ושיפורים, אבל לא תמיד תיעדתי בצורה מסודרת כל שינוי קטן
בהיפר־פרמטרים, בנתונים או באוגמנטציה. בדיעבד, תיעוד כזה היה מקל עליי
גם בכתיבת הדוח וגם בהשוואה בין גרסאות שונות של המערכת.

בנוסף, הייתי משקיע מוקדם יותר באיסוף של יותר נתונים "אמיתיים" מתוך
סרטוני תחרות. אחד הלקחים הגדולים שלי מן הפרויקט הוא שאיכות וריאליזם של
הנתונים חשובים לא פחות מבחירת המודל.

\subsection{מה למדתי על עצמי}

הפרויקט לימד אותי שאני נהנה במיוחד מעבודות שמחברות בין חשיבה
אלגוריתמית, פיתוח תוכנה וראייה ממוחשבת. גיליתי שאני אוהב לא רק ליישם
מודלים קיימים, אלא גם להבין לעומק את מגבלותיהם ולבנות סביבם פתרון
הנדסי רחב יותר.

הוא גם לימד אותי סבלנות. תיוג נתונים, בדיקות חוזרות וניסויים שלא עבדו
היו חלק משמעותי מהדרך, ודווקא הם חיזקו אצלי את התחושה שאני מסוגל
להמשיך לעבוד גם כשהפתרון אינו ברור מיידית.

\subsection{המשך הדרך}

הפרויקט חיזק אצלי מאוד את העניין בלמידת מכונה ובראייה ממוחשבת, אבל גם
בתחום הרחב יותר של מערכות חכמות. הוא גרם לי להבין שאני רוצה להמשיך
לעסוק בתחומים שבהם אלגוריתמים פוגשים בעיה אמיתית מן העולם, ולא נשארים
רק ברמת הניסוי התיאורטי.

אם אמשיך את הפרויקט בעתיד, הייתי רוצה להפוך אותו למערכת מלאה יותר:
עם הערכה כמותית מסודרת, שיפור מודל הספרות ואולי גם ניתוח מתקדם יותר של
התנהגות רובוטים על המגרש.
